二值化后统计连通域计数时，相互接触的米粒会被并作一粒，米粒越密漏计越多。
针对这一问题，本文提出一种以图像自身尺度为基准的粘连米粒计数方法。
该方法先由图像量出单粒米的面积与形状，据此判定哪些连通域为粘连块，并控制注水分割的种子数量。
切分之后再按面积与形状复核结果，并剔除硬币、边框等非米目标。
在真实照片、自制可控粘连合成图与低分辨率图像三组数据上，平均绝对误差分别为
4.73、1.77 与 1.93 粒，优于五种经典对比方法。
粘连率为 80\% 时，误差约为直接统计连通域的十一分之一。
此外，本文以视觉大模型 SAM 3 的输出作为伪标签，蒸馏出参数量 27 万的轻量网络。
它的总体误差低于上述几何方法，而模型体积约为 SAM 3 的三千分之一。
文中还记录了学习过程中经测试否定的若干设想，以及尚未解决的问题。
